% ============================================================================
% ARBEITSBLATT DS 5 - Übung ser/estar/hay (während Wartezeit)
% ============================================================================
% Ziel: 15 Punkte, 15 Minuten Bearbeitungszeit
% Für SuS, die auf ihre Prüfung warten
% ============================================================================

\documentclass[11pt,a4paper]{article}

\usepackage[utf8]{inputenc}
\usepackage[T1]{fontenc}
\usepackage[ngerman]{babel}
\usepackage[left=2cm,right=2cm,top=1.8cm,bottom=1.5cm]{geometry}
\usepackage{helvet}
\usepackage[table]{xcolor}
\usepackage{tabularx}
\usepackage{array}
\usepackage{enumitem}
\usepackage{tcolorbox}
\usepackage{fancyhdr}
\usepackage{lastpage}
\usepackage{amssymb}

\renewcommand{\familydefault}{\sfdefault}

\definecolor{spanisch}{HTML}{c41e3a}
\definecolor{grau}{HTML}{666666}
\definecolor{hellgrau}{HTML}{f5f5f5}

\pagestyle{fancy}
\fancyhf{}
\fancyhead[L]{\small\textcolor{grau}{Spanisch Spätbeginner -- Übung}}
\fancyhead[R]{\small\textcolor{grau}{AB-05}}
\fancyfoot[C]{\small\thepage/\pageref{LastPage}}
\renewcommand{\headrulewidth}{0.5pt}

\newcommand{\luecke}[1]{\underline{\hspace{#1}}}
\newcommand{\punktebox}[1]{\hfill\fbox{\small \_\_/#1}}
\newcommand{\es}[1]{\textit{#1}}

\newtcolorbox{infokasten}[1][]{
    colback=hellgrau,
    colframe=grau,
    fonttitle=\bfseries,
    title=#1,
    boxrule=0.5pt,
    arc=0pt,
    left=5pt, right=5pt, top=3pt, bottom=3pt
}

\newcounter{aufgabe}
\newcommand{\aufgabe}[2]{%
    \stepcounter{aufgabe}%
    \vspace{0.8em}%
    \noindent\textbf{\theaufgabe.} #1 \punktebox{#2}%
    \vspace{0.3em}%
}

\begin{document}

% ============================================================================
% KOPFBEREICH
% ============================================================================
\noindent
\begin{tabularx}{\textwidth}{@{}Xr@{}}
    {\Large\textbf{\textcolor{spanisch}{Repaso: ser -- estar -- hay}}} & Name: \luecke{5cm} \\[0.3em]
    {\small DS 5 -- Übung während Wartezeit (15 Min)} & Datum: \luecke{3cm} \\
\end{tabularx}

\vspace{0.3em}
\hrule
\vspace{0.8em}

\begin{infokasten}[Zur Erinnerung]
\begin{tabularx}{\textwidth}{@{}|p{2.5cm}|X|X|@{}}
\hline
\rowcolor{hellgrau}
\textbf{Verb} & \textbf{Verwendung} & \textbf{Beispiele} \\
\hline
\textbf{ser} & Eigenschaften, Identität & Mi barrio \textbf{es} tranquilo. \\
\hline
\textbf{estar} & Ort, Zustand & La iglesia \textbf{está} en el centro. \\
\hline
\textbf{hay} & Existenz (es gibt) & \textbf{Hay} un parque grande. \\
\hline
\end{tabularx}
\end{infokasten}

\vspace{0.8em}

% ============================================================================
% AUFGABE 1: Lückensätze (10 Sätze)
% ============================================================================
\aufgabe{Ergänze \es{ser}, \es{estar} oder \es{hay} in der richtigen Form.}{10}

\begin{enumerate}[label=\alph*), leftmargin=2em, itemsep=0.4em]
    \item Mi barrio \luecke{2cm} muy bonito y tranquilo.
    \item En mi calle \luecke{2cm} una panadería y un supermercado.
    \item La biblioteca \luecke{2cm} cerca de mi casa.
    \item ¿Cómo \luecke{2cm} tu barrio? -- Mi barrio \luecke{2cm} moderno.
    \item El parque \luecke{2cm} al lado del cine.
    \item No \luecke{2cm} piscina en mi barrio.
    \item La farmacia \luecke{2cm} abierta hasta las 20:00.
    \item ¿Dónde \luecke{2cm} el hospital? -- \luecke{2cm} en la calle Mayor.
    \item Las tiendas \luecke{2cm} pequeñas pero muy buenas.
    \item ¿\luecke{2cm} un banco cerca de aquí? -- Sí, \luecke{2cm} uno en la plaza.
\end{enumerate}

\vspace{1em}

% ============================================================================
% AUFGABE 2: Zuordnung Orte + Beschreibungen
% ============================================================================
\aufgabe{Zuordnung: Verbinde die Orte mit passenden Beschreibungen.}{5}

\vspace{0.3em}

\begin{tabularx}{\textwidth}{|p{4cm}|p{0.8cm}|p{9cm}|}
\hline
\rowcolor{hellgrau}
\textbf{Ort} & \textbf{Nr.} & \textbf{Beschreibung} \\
\hline
1. el supermercado & \luecke{0.5cm} & A. Aquí hay libros y es tranquilo. \\
\hline
2. el parque & \luecke{0.5cm} & B. Aquí hay médicos y enfermeras. \\
\hline
3. la biblioteca & \luecke{0.5cm} & C. Aquí hay árboles y bancos para sentarse. \\
\hline
4. el hospital & \luecke{0.5cm} & D. Aquí hay frutas, verduras y pan. \\
\hline
5. la piscina & \luecke{0.5cm} & E. Aquí hay agua y puedes nadar. \\
\hline
\end{tabularx}

\vspace{1.5em}
\hrule
\vspace{0.5em}

\begin{flushright}
\textbf{Gesamt: \luecke{1cm} / 15 Punkte}
\end{flushright}

\vspace{0.5em}

\begin{infokasten}[Tipp für die Prüfung]
\begin{itemize}[leftmargin=1.5em, nosep]
    \item \textbf{ser:} ¿Cómo \textbf{es}...? (Eigenschaften beschreiben)
    \item \textbf{estar:} ¿Dónde \textbf{está}...? (Ort angeben)
    \item \textbf{hay:} ¿\textbf{Hay}...? (Existenz fragen)
\end{itemize}
\end{infokasten}

\end{document}
